\section{Threats to Validity}
\label{sec:threats}

In this section, we discuss potential threats to the validity of our empirical study and how they were mitigated.

\subsection{Internal Validity}
Threats to internal validity concern factors that could have influenced our experimental results without our knowledge. In LLM-based generation, the non-deterministic nature of language models poses a threat. We mitigated this by setting the generation temperature parameter to $0.0$ and disabling sampling, ensuring that the model output is deterministic and reproducible.

\subsection{External Validity}
Threats to external validity concern the generalizability of our findings. We evaluated Qwen2.5-7B-Instruct on $N=100$ user stories from the Rathnayake dataset~\cite{rathnayake2026}. While these requirements represent a diverse set of web application scenarios, the results may not generalize to other software engineering domains (such as safety-critical embedded systems or real-time gaming engines) or other programming languages (such as Java or C++).

\subsection{Construct Validity}
Threats to construct validity concern whether our metrics actually measure what they are intended to measure. We used SBERT Cosine Similarity as an automated proxy for semantic alignment, which may not capture all business logic nuances compared to human inspection. We mitigated this by manually validating a subset of the pilot samples and confirming that SBERT scores $\ge 0.80$ align well with human judgments. For syntax checking, AST parsing verifies syntactic correctness but does not check if the Python step definitions are functionally complete (e.g., whether they contain actual web driver locators instead of empty \texttt{pass} statements).

\subsection{Conclusion Validity}
Threats to conclusion validity concern the statistical power of our tests. We mitigated this by selecting $N=100$ samples, which is statistically sufficient to run non-parametric tests (such as Wilcoxon and McNemar tests). We also reported p-values and effect sizes (Cohen's d) to ensure that the magnitude of difference is clearly quantified.
